\par\noindent\textbf{(1)} 取仿射参数 \(s\)，测地线Lagrangian为 \(\mathcal L=\tfrac{R^2}{2}(\theta'^2+\sin^2\theta\,\phi'^2)\)。变分给 \(\theta''-\sin\theta\cos\theta\,\phi'^2=0\) 与 \((\sin^2\theta\,\phi')'=0\)。弧长归一时 \(R^2(\theta'^2+\sin^2\theta\,\phi'^2)=1\)。
\par\noindent\textbf{(2)} 因 \(\phi\) 为循环坐标，\(p_\phi=R^2\sin^2\theta\,\phi'\) 恒定。若 \(\alpha\) 为曲线切线与纬线方向的夹角，沿纬线的弧长分量为 \(R\sin\theta\,\phi'=\cos\alpha\)，故 \(\sin\theta\cos\alpha=p_\phi/R\) 为常量，这就是球面Clairaut关系。
\par\noindent\textbf{(3)} 任意经过球心的平面与球面相交为大圆。选择平面为赤道时 \(\theta=\pi/2,\ \phi=s/R+\phi_0\)，它满足两条测地线方程；由于球面旋转对称，任意大圆由赤道旋转得到，也必满足测地线方程。小圆纬线一般不是自由测地线。
\par\noindent\textbf{(4)} 球面的高斯曲率 \(K=1/R^2\)。沿单位速测地线的垂直Jacobi场幅度 \(j(s)\) 满足 \(j''+Kj=0\)，若 \(j(0)=0,j'(0)\ne0\)，则 \(j(s)\propto\sin(s/R)\)。首个共轭点出现在 \(s=\pi R\)，即反足点处，大圆局部最短性质在此之后失去全局唯一性。
